%packages
\usepackage{amsmath, amsthm, amssymb}
\usepackage{graphicx}
\usepackage{enumerate}
\usepackage[utf8]{inputenc}
% \usepackage[spaced=0]{diffcoeff}
\usepackage[italic=true]{derivative}
\usepackage{tcolorbox}
\usepackage{physics2}
\usephysicsmodule{ab, braket}
\usepackage{siunitx}
\usepackage{fancyhdr}
\usepackage[margin=1in]{geometry}
\usepackage{tikz}
\usepackage{pgfplots} % for the axis environment
% \usepackage{fullpage}

%environments
\newtheorem*{question}{Q}
\newtheorem*{prop}{Proposition}

\newenvironment{quest}[1]
{\vspace{0.3cm}\noindent\textbf{Q(#1)}\em}{\vspace{0.3cm}}

\newenvironment*{Remark}{\noindent\textbf{Remarks}\begin{enumerate}}{\end{enumerate}\vspace{1em}\noindent}


\usepackage{hyperref}
\hypersetup{
    colorlinks=true,
    linkcolor=blue,
    filecolor=magenta,      
    urlcolor=cyan,
    pdftitle={Overleaf Example},
    pdfpagemode=FullScreen,
}


%commands
\newcommand{\e}{\varepsilon}
\newcommand{\R}{\mathbb{R}}
\newcommand{\Z}{\mathbb{Z}}
\newcommand{\N}{\mathbb{N}}
\newcommand{\C}{\mathbb{C}}
\newcommand{\Q}{\mathbb{Q}}


\newcommand{\grad}{\vec \nabla}
\newcommand{\LL}{\mathcal{L}}
\newcommand{\vphi}{\varphi}
\newcommand{\p}{\partial}



\newcommand{\eq}{\Leftrightarrow}
\newcommand{\clm}{\textbf{Claim }}
\newcommand{\rmks}{\noindent\textbf{Remarks}}

\newcommand{\levi}{\mathfrak{E}}


\makeatletter
     \newcommand\vb{\@ifstar\boldsymbol\mathbf}
     \newcommand\va[1]{\@ifstar{\vec{#1}}{\vec{\mathrm{#1}}}}
     \newcommand\vu[1]{%
       \@ifstar{\hat{{#1}}}{\hat{{#1}}}}
    % \renewcommand*{\vec}[][default]{def}
\makeatother


\newcommand{\ev}[1]{\langle{#1}\rangle}
\newcommand{\vecp}[1]{\vec{#1}\:'}
%\renewcommand\qedsymbol{QED}

%%%%%%%%%%%%%%%%%%%%%%
% Set up header/footer
\pagestyle{fancy}
\fancyhead[LO,L]{Vishwanath Wimalasena}
\setlength{\headheight}{14.5pt}
\fancyfoot[LO,L]{}
\fancyfoot[CO,C]{\thepage}
\fancyfoot[RO,R]{}
\renewcommand{\headrulewidth}{0.5pt}
\renewcommand{\footrulewidth}{0.4pt}
\renewcommand{\headrulewidth}{0.5pt}
%%%%%%%%%%%%%%%%%%%%%%



\tcbuselibrary{theorems}
\tcbuselibrary{breakable}
\tcbuselibrary{skins}

%environments
\newtheorem{theorem}{Theorem}[section]
\newtheorem{proposition}{Proposition}[section]
\newtheorem{definition}{Definition}[section]
\newtheorem{lemma}{Lemma}[section]
\newtheorem{lem}{Lemma}
\newtheorem*{ex}{Exercise}

\theoremstyle{definition}
\newtheorem*{sol}{Sol}

\newtcbtheorem[number within=section]{cthm}{Theorem}%
{colback=green!8!white, colframe=green!40!black,fonttitle=\bfseries, fontupper=\itshape}{mainthm}

\newtcbtheorem[number within=section]{cexercise}{Exercise}%
{colback=green!8!white, colframe=green!40!black,fonttitle=\bfseries, fontupper=\itshape}{mainthm}

\newtcbtheorem[number within=section]{cclaim}{Claim}%
{colback=green!8!white,colframe=green!40!black,fonttitle=\bfseries, fontupper=\itshape}{claim}

\tcolorboxenvironment{sol}{breakable, enhanced, left=5mm,
    before skip=10pt, after skip=10pt, frame  hidden,colback=green!6!white,
    borderline west={1pt}{0pt}{green!60!black}}

\tcolorboxenvironment{lem}{
    enhanced jigsaw,colback=green!8!white, colframe=green!40!black,breakable,before skip=10pt,after skip=10pt }

\tcolorboxenvironment{proof}{% `proof' from `amsthm' 
    blanker,breakable,left=5mm,
    before skip=10pt,after skip=10pt,
    borderline west={1pt}{0pt}{green!60!black}}

% \newtcolorbox{ansbox}{tcbox raise base, leftrule=0pt, arc=0pt, rightrule=0pt,toprule=0mm,bottomrule=0mm, colback=green!8!white, colframe=green!20!black}

% \newtcolorbox{background}[1][]{enhanced,
%     colback=green!5!white,colframe=green!40!black,
%     colbacktitle=green!40!black,fonttitle=\bfseries,
%     frame hidden,
%     title=Background,
%     boxrule=0pt,
%     titlerule style=green!70!black, #1}

\tcolorboxenvironment{ex}{enhanced jigsaw, colback=green!8!white, colframe=green!40!black,breakable,before skip=10pt,after skip=10pt}

\tcolorboxenvironment{question}{% `proof' from `amsthm' 
    blanker,colback=green,breakable,left=5mm,
    before skip=10pt,after skip=10pt,
    borderline west={1mm}{0pt}{green!40!black}}

% \tcbset{one-sided-color/.style={
%         breakable,
%         enhanced,
%         boxrule=0pt,
%         frame hidden,
%         borderline ={1pt}{0pt}{#1},
%         colback=#1!10,
%         coltitle=#1!200,
%         after title={\ },
%         attach title  to upper,
%     }
% }

\newtcolorbox{q}{one-sided-color=green!60!black, title=\textbf{Q}}

\newtcolorbox{outline}[1]{one-sided-color=blue!70!black, title=\textbf{#1}}

% \newtcolorbox{ans}{breakable,
%         enhanced,
%         boxrule=0pt,
%         frame hidden,
%         borderline west={1pt}{0pt}{green!70!black},
%         % borderline south={1pt}{0pt}{green!70!black},
%         % after title={\ },
%         % attach title  to upper,
%         top=1pt,
%         % boxsep=1pt,
%         colback=white,
%         coltitle=green!70!black!200}


% \newtcolorbox{q2}{breakable,
%         enhanced,
%         boxrule=0pt,
%         frame hidden,
%         borderline west={1pt}{0pt}{green!70!black},
%         colback=green!70!black!10, after title={Q}}

% \newtcolorbox{exc}[1]{one-sided-color=green!70!black, title=\textbf{(#1)}}
\definecolor{bordercolor}{rgb}{0.13,0.57,0.07}
\definecolor{boxcolor}{rgb}{0.28,0.69,0.28}

\newtcolorbox{exc}[1]{
        breakable,
        enhanced,
        frame hidden,
        borderline ={1pt}{0pt}{bordercolor},
        colback=boxcolor!10,
        coltitle=bordercolor,
        after title={\ },
        attach title  to upper,
        title=\textbf{Exercise #1},
        }

\newtcolorbox{ans}{
    breakable,
    enhanced,
    boxrule=0pt,
    frame hidden,
    borderline west={1pt}{0pt}{bordercolor},
    top=0pt,
    colback=white,
    coltitle=bordercolor,
    }
% \newcommand{\exc}[1]{\begin{exc}#1\end{exc}}